\section{Related Work and Literature Review}
\label{sec:related_work}
To establish the theoretical and algorithmic foundations of ROUTE/IQ-NextGen, we systematically reviewed 27 pivotal research papers across four foundational domains, detailed in Table~\ref{tab:lit_review_27}.

\subsection{Classical Vehicle Routing and Heuristic Methods}
The foundations of combinatorial vehicle routing stem from the classic formulations of Toth and Vigo~\cite{toth2014vehicle}. Clarke and Wright~\cite{clarke1964scheduling} established the greedy edge-savings heuristic, which remains a standard baseline for rapid initialization. Croes~\cite{croes1958method} developed the foundational 2-opt edge exchange for the Traveling Salesperson Problem, which forms the basis of route edge improvements. Kuhn~\cite{kuhn1955hungarian} provided the exact Hungarian primal-dual method for bipartite assignment. Gillett and Miller~\cite{gillett1974sweep} introduced the polar coordinate sweep clustering algorithm, enabling geometric partitioning of delivery stops around a central depot. 

To overcome greedy local entrapment, Glover~\cite{glover1986future} formulated Tabu Search, introducing recency and frequency memory vectors to prevent cyclic revisiting of candidate states. Voudouris and Tsang~\cite{voudouris1999guided} proposed Guided Local Search (GLS), which augments the objective function with penalty terms to guide local search out of local minima. Hansen and Mladenovi{\'c}~\cite{hansen2001variable} unified multi-neighborhood search via Variable Neighborhood Search (VNS) and Variable Neighborhood Descent (VND). Finally, the Bureau of Public Roads (BPR)~\cite{bpr1964traffic} formulated the empirical non-linear volume-delay function governing urban congestion.

\subsection{Big Data and Spatial Processing Architectures}
Modern urban routing operates on massive order volumes requiring distributed processing. Zaharia et al.~\cite{zaharia2016apache} introduced Apache Spark, utilizing in-memory Resilient Distributed Datasets (RDDs) and lineage graphs to achieve fault-tolerant parallel computation. Yu et al.~\cite{yu2015geospark} extended Spark into GeoSpark (now Apache Sedona), integrating spatial RDDs and spatial partitioning trees for large-scale geospatial range and join queries.

\subsection{Quantum Computing and Variational Algorithms in NISQ}
Variational quantum algorithms have emerged as candidates for combinatorial optimization on Noisy Intermediate-Scale Quantum (NISQ) devices~\cite{preskill2018quantum}. Farhi, Goldstone, and Gutmann~\cite{farhi2014qaoa} introduced the Quantum Approximate Optimization Algorithm (QAOA), using alternating problem and mixer Hamiltonians. Hadfield et al.~\cite{hadfield2019qaoa} extended QAOA into the Quantum Alternating Operator Ansatz, designing constrained $XY$-mixers. Wang et al.~\cite{wang2020xy} analyzed hard versus soft constraint handling with $XY$-mixers. 

Feld et al.~\cite{feld2019hybrid} and Borowski et al.~\cite{borowski2020new} explored hybrid quantum-classical decomposition for CVRP using quantum annealers. Harwood et al.~\cite{harwood2021formulating} analyzed QAOA constraint penalties. Azad et al.~\cite{azad2022solving} applied QAOA to small symmetric VRP instances ($N \le 8$). Glos et al.~\cite{glos2022space} introduced space-efficient binary encodings for variational circuits. Lucas~\cite{lucas2014ising} mapped canonical NP-complete problems to Ising spin glass models. However, McClean et al.~\cite{mcclean2018barren} proved the emergence of barren plateaus in parameterized quantum circuits, highlighting training challenges. Qiskit~\cite{qiskit2024} provides the open-source infrastructure for simulating and executing these quantum algorithms.

\subsection{Cooperative Search, Hyper-Heuristics and Adaptive Learning}
Cooperative search was formalized by Toulouse and Crainic~\cite{toulouse2004cooperative}, defining paradigms where autonomous heuristic agents exchange subproblems and solutions via central blackboards. Talbi~\cite{talbi2009metaheuristics} provided a comprehensive taxonomy of hybrid metaheuristics. Burke et al.~\cite{burke2013hyper} surveyed hyper-heuristic frameworks that select or generate heuristics automatically. Auer, Cesa-Bianchi, and Fischer~\cite{auer2002finite} proved finite-time logarithmic regret bounds for the Upper Confidence Bound (UCB-1) multi-armed bandit algorithm, which provides the mathematical foundation for our adaptive operator controller.
